\chapter{Arquitectura Multi-Tenant}

\section{Descripci\'on General}

Vetrinaria implementa una arquitectura multi-tenant basada en filas
(row-based), donde todas las cl\'{\i}nicas comparten las mismas tablas
pero los datos se a\'{\i}sla mediante la columna \texttt{clinic\_id}
en todas las entidades.

\section{Estrategia de Aislamiento}

\subsection{Tenancy por Filas}

Cada tabla de datos tiene una columna \texttt{clinic\_id} que act\'ua
como foreign key hacia la tabla \texttt{clinics}. Las consultas
cr\'{\i}ticas filtran por \texttt{clinic\_id} usando el contexto de la
sesi\'on del usuario autenticado.

\begin{longtable}{|l|l|l|}
\hline
\textbf{Tabla} & \textbf{Clinic-scoped} & \textbf{Notas} \\
\hline
clinics & N/A & Tabla ra\'{\i}z \\
users & S\'{\i} & \\
clients & S\'{\i} & \\
patients & S\'{\i} & \\
appointments & S\'{\i} & \\
medical\_records & S\'{\i} & \\
dental\_records & S\'{\i} & \\
vaccines & S\'{\i} & \\
vaccination\_records & S\'{\i} & \\
products & S\'{\i} & \\
suppliers & S\'{\i} & \\
invoices & S\'{\i} & \\
invoice\_items & No (hereda de invoice) & \\
notifications & No (por usuario) & \\
\hline
\end{longtable}

\section{Contexto Multi-Tenant en la Sesion}

El ID de la cl\'{\i}nica se inyecta en el JWT durante la autenticaci\'on:

\begin{lstlisting}[style=typescript, caption=JWT callback - Inyeccion de clinicId]
async jwt({ token, user }) {
  if (user) {
    (token as any).role = user.role;
    (token as any).id = user.id;
    (token as any).clinicId = user.clinicId;
  }
  return token;
},
\end{lstlisting}

La utilidad \texttt{getCurrentClinicId()} extrae este valor:

\begin{lstlisting}[style=typescript, caption=Obtencion del clinicId desde la sesion]
export async function getCurrentClinicId(): Promise<number> {
  const session = await auth();
  if (!session?.user?.clinicId) {
    throw new Error('No clinic context available');
  }
  return session.user.clinicId;
}
\end{lstlisting}

\section{Filtrado en Consultas}

Todas las server actions que acceden a datos multi-tenant siguen este
patr\'on:

\begin{lstlisting}[style=typescript, caption=Patron de filtrado por clinica]
export async function getClients(search?: string) {
  const clinicId = await getCurrentClinicId();
  const clinicFilter = eq(schema.clients.clinicId, clinicId);

  return db
    .select()
    .from(schema.clients)
    .where(clinicFilter)
    .orderBy(schema.clients.name);
}
\end{lstlisting}

\section{Rol Super Admin}

El rol \texttt{super\_admin} tiene acceso global a todas las cl\'{\i}nicas
y puede gestionar:

\begin{itemize}[leftmargin=*]
    \item Creaci\'on y edici\'on de cl\'{\i}nicas.
    \item Acceso a datos de cualquier cl\'{\i}nica.
    \item Configuraci\'on global del sistema.
\end{itemize}

\section{Migraci\'on Multi-Tenant}

La migraci\'on \texttt{0003\_multi-tenant.sql} realiz\'o:

\begin{enumerate}
    \item Creaci\'on de la tabla \texttt{clinics}.
    \item Inserci\'on de la cl\'{\i}nica por defecto ``Cl\'{\i}nica Principal''.
    \item Adici\'on de \texttt{clinic\_id} nullable a 8 tablas.
    \item Actualizaci\'on de registros existentes con el ID por defecto.
    \item Establecimiento de NOT NULL y llaves for\'aneas.
\end{enumerate}

\section{Gesti\'on de Cl\'inicas}

La p\'agina \texttt{/settings/clinics} permite a super admins gestionar
cl\'{\i}nicas:

\begin{itemize}[leftmargin=*]
    \item Lista de cl\'{\i}nicas con nombre, slug, email.
    \item Creaci\'on con validaci\'on de slug (\texttt{[a-z0-9-]+}).
    \item Edici\'on de datos y configuraci\'on.
\end{itemize}

\subsection{Server Actions}

\begin{lstlisting}[style=typescript, caption=Clinicas - Server actions con Zod]
const clinicSchema = z.object({
  name: z.string().min(1, 'El nombre es obligatorio'),
  slug: z.string().min(1, 'El slug es obligatorio')
    .regex(/^[a-z0-9-]+$/,
           'Solo minusculas, numeros y guiones'),
  email: z.string().email().optional().or(z.literal('')),
  phone: z.string().optional(),
  address: z.string().optional(),
});
\end{lstlisting}
